
\subsubsection{Software Ecosystem Dynamics}

Studies of software package ecosystems provide conceptual foundations for lifecycle analysis. Wittern et al. (2016) analyzed npm package dynamics, identifying distinct lifecycle patterns in package downloads. Decan et al. (2019) conducted comparative analysis across seven packaging ecosystems, establishing that adoption trajectories vary with ecosystem characteristics. Mujahid et al. (2021) identified breaking update patterns through lifecycle analysis of npm packages. These studies establish precedent for lifecycle-based ecosystem analysis but focus on software packages rather than other time series domains.

\subsubsection{Time Series Clustering}

Time series clustering methods have evolved from distance-based to shape-aware approaches. Aghabozorgi et al. (2015) provide a comprehensive review distinguishing between raw-data, feature-based, and model-based clustering methods. DTW-based clustering, formalized by Berndt and Clifford (1994), has become a standard approach for shape-aware time series analysis. The tslearn library \citep{Tavenard2020tslearn} provides implementations that enable reproducible time series clustering research.

\subsubsection{Changepoint Detection}

Changepoint detection identifies points where statistical properties of time series shift. The PELT algorithm \citep{Killick2012Optimal} achieves optimal detection with linear computational cost. PELT has been applied to climate analysis, financial time series, and network traffic patterns.

The contribution of this work is methodological integration: combining PELT changepoint detection with DTW clustering in a two-level hierarchy that separates phase detection from trajectory classification.
